\section{Research objective, source baseline, and quantitative target}
\label{sec:context}

The immediate objective is to turn compressed normal-surface data into more
useful, independently checked geometric operations. The eventual objective
is a deterministic unknot recognizer with a proved quasi-polynomial
worst-case bound in the number of crossings of its input diagram. These
objectives are connected, but they require different theorems. A fast orbit
routine, a successful terminal handlebody test, or a polynomial-size positive
certificate does not by itself control the total number of surfaces or
hierarchy states an algorithm must inspect.

This contribution supplies a new source-relative compiler step and a
separate, exact certificate of a sufficient condition for hierarchy progress.
The first is integrated into the existing cut-classification pipeline. The
second is an algebraic research component whose future geometric interface
is specified explicitly. Their proofs, implementation, negative examples,
and experimental limits are part of one deliverable.

\subsection{The pinned state of ProveIt}

All baseline comparisons use repository revision
\begin{center}
\texttt{8cb7d9d25ac8cb64a9d147304e3441e98e94c7a4}.
\end{center}
The snapshot is from 12 October 2026. The 905 original runtime-directory
blobs were authenticated against their Git
blob identifiers. The delivered patch is relative to that revision, rather
than an unpinned moving branch. The incoming research packages in the same
snapshot were also inspected before choosing the present contribution
\cite{repo,prior-genus,prior-stopping,prior-saturated,prior-hierarchy,prior-sparse}.

The relevant existing pipeline starts from a compact, connected, orientable
finite triangulation with one torus boundary and an admissible normal
coordinate vector. It constructs the actual cut in compressed form:
a bounded complementary handle core, interval-bundle families with binary
multiplicities, and authenticated ordered interfaces. The terminal operator
then certifies a census of handlebody components when its geometric
cancellation and bundle-attachment criteria apply.

The word \Complete\ in that operator means an exact classification of the
\emph{supplied cut}: every component is certified to be a handlebody and its
genus and multiplicity are computed. It is not a complete decision procedure
for arbitrary handlebodies, and it is not a standalone knot verdict. A
nontrivial knot exterior can have a cut consisting of handlebodies. The word
\Inconclusive\ records the failure of this particular positive compiler,
not a proof that a component is a non-handlebody or that a knot is nontrivial.

The pinned terminal operator certifies 85 of the maintained 106 source cuts.
Ten unresolved inputs are rejected immediately because the source core
contains three-handles. The first theorem family below removes that entire
\emph{structural refusal}: every source three-handle can be cancelled after
the normal cut. Later positive cancellation still need not complete. On this
corpus seven of the ten inputs become certified, giving 92 successes.

\subsection{Prior results that this article uses}

The following distinctions prevent repeated work from being presented as a
new discovery.

\begin{itemize}
\item The incoming \emph{Genus Budgets and Native Interfaces} contribution
already proves completeness of the ordinary source two/three queue before
a normal cut and with the initial empty boundary pattern. Classical
two/three cancellation is also prior work. The new result concerns the
\emph{compressed complementary core after an arbitrary normal cut}, its
endpoint-gap side table, preservation of bundle ports, and independent
source replay.
\item The unsigned stopping-set characterization of positive raw one/two
cancellation, complete column peeling, and the associated unit-permanent
minor criterion are already in the incoming stopping-core work. Our new
use couples those obstructions to the set of allowable source dual forests
and realizes a strict forest-choice advantage in an actual small manifold.
\item Saturated bundle cancellation is already incorporated in the pinned
baseline. Consequently, 78 successes is an obsolete comparison point for
this patch. The appropriate baseline is 85.
\item The chain-space separation theorem used for the second component is
finite-dimensional rational linear algebra. The contribution is its exact
certificate contract, independent integer replay without trusting a
purported complete bundle-cycle basis, and its precise use with Lackenby's
equality case. We do not claim a new general linear-algebra duality theorem.
\end{itemize}

These results extend the repository's current research programme. Assertions
of priority beyond that programme would require a broader historical review;
none is needed for the geometric and algorithmic claims proved here.

\subsection{What exactly has Lackenby announced?}
\label{subsec:announcements}

The available primary sources contain two quantitative formulations. The
February 2021 talk displays a bound of the form
\[
 n^{O(\log n)}=\exp\bigl(O((\log n)^2)\bigr),
\]
and discusses logarithmic hierarchy depth. The March 2021 Oxford handout
instead displays
\[
 2^{O((\log n)^3)}=n^{O((\log n)^2)},
\]
with a corresponding squared-logarithmic depth formulation. Here $n$ is the
crossing number of the supplied diagram. Both bounds are quasi-polynomial,
but they are quantitatively different; the February bound is the stronger
target \cite{lackenby2021feb,lackenby2021mar}. The Oxford announcement also
describes an $n^{c\log n}$ bound \cite{oxford2021}.

The July 2026 preprint \emph{Incompressible surfaces, hierarchies and unknot
recognition}, already included among the project's source material, develops
the geometric hierarchy machinery \cite{lackenby2026}. Its Section 9
explicitly leaves running-time estimates outside the precision of the stated
algorithm. In particular, a quantitative iteration bound and a linear bound
in quadratic handle complexity are not themselves either 2021 asymptotic
claim. This article uses the preprint's stated theorems, while treating the
global quasi-polynomial estimate as a research target requiring additional
proofs.

\subsection{Parameters and levels of certification}

We distinguish four sizes throughout the article:
\[
\begin{aligned}
 n&=\text{diagram crossings},&
 t&=\text{source tetrahedra},\\
 L&=\text{maximum coordinate bit length},&
 W&=\sum_i x_i.
\end{aligned}
\]
The source input has size $O(t\log(t+1)+tL)$ in a direct encoding. Normal
weight $W$ may be exponential or larger relative to $t$; a compressed
operation must charge arithmetic to $L$, not loop once per normal disc.
Later chain-matrix results use their explicit ranks and entry bit lengths
as separate input parameters. There is no established identification of all
those ranks with a polynomial in $n$ along an arbitrary future hierarchy.

Three logical levels must also remain distinct. A chart-level verifier
authenticates geometric incidence from the triangulation and coordinates.
A local certificate, such as a dual forest, proves a transformation of that
authenticated geometry. A global recognition theorem must additionally
establish that the algorithm finds an adequate sequence of such operations,
terminates correctly on both answers, and spends the claimed total time.
The delivered native patch implements the first two levels for a new useful
transformation. The final sections specify the missing work at the third.
